\chapter{The Born Hierarchy}
\label{ch:born-rule}

Up to this point, this book has reconstructed \emph{constraints}: what
admissibility permits and forbids, without reference to the specific
formulas a physics textbook opens with. From here onward, the project
reverses direction — using those constraints to reconstruct the familiar
formalism itself, formula by formula, with each one's exact cost made
visible rather than assumed. Chapter~\ref{ch:assumption-audit} flagged four distinct statements this book
has been careful not to conflate. This chapter is not "the Born-rule
derivation"; it is a dependency chain, with each link's assumptions made
explicit rather than absorbed silently into the next.

\section{The General Context-Transition Identity}

Chapters~\ref{ch:context-graphs}--\ref{ch:real-mub-obstruction} established
the context-transition interpretation of $B=|U|^2$ only for real
two-dimensional frames (rotations in $\R^2$). The complex, general-dimension
statement is an easy generalization, but it has not yet been written down,
and Chapter~\ref{ch:born-rule} depends on it, so it is stated explicitly here
rather than assumed.

\begin{proposition}[Context-transition identity, general $d$]
\label{prop:general-context-transition}
\index{Context-transition identity, general d (proposition)}
Let $F_A = (|a_1\rangle,\dots,|a_d\rangle)$ and $F_B =
(|b_1\rangle,\dots,|b_d\rangle)$ be orthonormal bases of $\C^d$. Then $U :=
F_A^\dagger F_B$, with entries $U_{ij} = \langle a_i|b_j\rangle$, is unitary,
and $B_{ij} := |\langle a_i|b_j\rangle|^2$ is unistochastic.
\end{proposition}

\begin{proof}
$(F_A^\dagger F_B)^\dagger(F_A^\dagger F_B) = F_B^\dagger F_A F_A^\dagger F_B
= F_B^\dagger F_B = I$, using unitarity of $F_A$ (so $F_AF_A^\dagger=I$) and
of $F_B$; hence $U$ is unitary. $B_{ij}=|U_{ij}|^2$ is unistochastic by
definition (Chapter~\ref{ch:unistochastic}).
\end{proof}

\section{From Contexts to Pure States}

Proposition~\ref{prop:general-context-transition} is a statement about
\emph{whole frames}. It says nothing yet about a single normalized vector
$|\psi\rangle$ carrying its own probabilities. Bridging this gap requires an
explicit operational assumption, stated here for the first time in the book:

\begin{assumption}[Pure states as frame-completable vectors]
\label{ass:frame-completion}
A prepared pure state $|\psi\rangle$ (normalized, $\langle\psi|\psi\rangle=1$)
can be operationally represented as the first vector of \emph{some}
orthonormal frame $F_\psi = (|\psi\rangle,|\psi_2\rangle,\dots,|\psi_d\rangle)$
— i.e.\ preparing $|\psi\rangle$ is representable as selecting one designated
column of an admissible context.
\end{assumption}

This is a genuine addition to the book's assumption base, not a consequence
of Parts I--II: nothing proved so far distinguishes "a state" from "a
context" as separate kinds of object, and this chapter's first substantive
move is to identify the former with a distinguished coordinate of the latter.
Any vector can be completed to an orthonormal basis (Gram--Schmidt), so the
mathematical content of Assumption~\ref{ass:frame-completion} is only the
operational claim that preparation is representable this way — worth keeping
visible precisely because it is doing real work.

\begin{theorem}[Pure-state Born rule]
\label{thm:pure-state-born}
\index{Pure-state Born rule (theorem)}
Given Assumption~\ref{ass:frame-completion}, and a measurement context $F_E =
(|e_1\rangle,\dots,|e_d\rangle)$,
\begin{equation}
p_i := B_{i1} = |\langle e_i|\psi\rangle|^2,
\end{equation}
where $B_{i1}$ is the first column of the context-transition matrix between
$F_E$ and $F_\psi$.
\end{theorem}

\begin{proof}
Apply Proposition~\ref{prop:general-context-transition} with $A=E$, $B=\psi$:
$U = F_E^\dagger F_\psi$ is unitary, with first column $U_{i1} =
\langle e_i|\psi\rangle$ (the inner product of $|e_i\rangle$ with the first,
designated column $|\psi\rangle$ of $F_\psi$). Then $B_{i1} = |U_{i1}|^2 =
|\langle e_i|\psi\rangle|^2$ directly from the definition of the
context-transition matrix.
\end{proof}

\begin{remark}
This is the intended payoff: the pure-state Born rule is not an independent
postulate about state vectors. It is literally one column of the
already-established context-transition identity, once a state is identified
with a frame via Assumption~\ref{ass:frame-completion}. The rest of $F_\psi$
(the vectors $|\psi_2\rangle,\dots$) plays no further role and can be
discarded; only the assumption that such a completion operationally
\emph{exists} was needed.
\end{remark}

\section{From Pure States to Density Operators}

Part II already established convex mixtures of pure states as legitimate
composite-system-independent objects (Chapter~\ref{ch:positivity}). For a
classical mixture $\rho = \sum_a q_a |\psi_a\rangle\langle\psi_a|$ ($q_a\ge0$,
$\sum_a q_a=1$), affine consistency under classical mixing — the probability
of outcome $i$ for a mixture is the corresponding mixture of probabilities —
gives
\begin{equation}
p_i = \sum_a q_a\, |\langle e_i|\psi_a\rangle|^2.
\end{equation}

\begin{theorem}[Density-operator Born rule]
\label{thm:density-born}
\index{Density-operator Born rule (theorem)}
$p_i = \operatorname{tr}(\rho P_i)$, where $P_i = |e_i\rangle\langle e_i|$.
\end{theorem}

\begin{proof}
$|\langle e_i|\psi_a\rangle|^2 = \langle\psi_a|e_i\rangle\langle
e_i|\psi_a\rangle = \operatorname{tr}\big(|\psi_a\rangle\langle\psi_a|\,P_i
\big)$ (cyclicity of trace). Summing with weights $q_a$ and using linearity of
trace, $p_i = \operatorname{tr}\Big(\big(\sum_a q_a
|\psi_a\rangle\langle\psi_a|\big) P_i\Big) = \operatorname{tr}(\rho P_i)$.
\end{proof}

\begin{remark}
No new assumption was needed for this step beyond
Theorem~\ref{thm:pure-state-born} and the affine-mixture structure Part II
already earned. This is the cleanest link in the chain.
\end{remark}

\section{What Is Not Yet Earned: POVMs}

\begin{equation}
p_i = \operatorname{tr}(\rho E_i), \qquad E_i \succeq 0,\ \ \sum_i E_i = I
\end{equation}
is a strictly more general statement than Theorem~\ref{thm:density-born}: it
permits $E_i$ that are not rank-one projectors, and in particular need not
come from any single orthonormal context at all. Nothing proved in this
chapter licenses that generality. The standard route — representing a
generalized measurement as an ordinary projective measurement on a larger,
composite system (a dilation) and invoking the partial trace of
Chapter~\ref{ch:partial-trace} — requires its own argument and is deferred
rather than assumed here.

\begin{ledgerentry}[End of Chapter~\ref{ch:born-rule}]
\textbf{Established:} a fully explicit dependency chain,
\begin{equation}
\boxed{|U_{ij}|^2 \ (\text{Ch.}\ref{ch:unistochastic}) \ \longrightarrow\
|\langle e_i|\psi\rangle|^2 \ (\text{Thm.}\ref{thm:pure-state-born}) \
\longrightarrow\ \operatorname{tr}(\rho P_i) \ (\text{Thm.}\ref{thm:density-born})}
\end{equation}
with each arrow's cost stated: the first arrow costs exactly Assumption
\ref{ass:frame-completion} (pure states as frame-completable vectors, a
genuinely new addition to the book's assumption base); the second arrow costs
nothing beyond what Part II already proved.\\
\textbf{Not established:} the POVM generalization $p_i =
\operatorname{tr}(\rho E_i)$, deferred to a later chapter pending a dilation
argument; the general complex Hilbert-space state structure for arbitrary $n$
(this chapter's propositions are dimension-generic but the book's worked
examples remain concentrated at $d=2,3$); Hermitian observables as the
general observable class, and spectral decomposition, remain untouched by
this chapter and are next.
\end{ledgerentry}
